\section{A Public Sequential-Cut Ceremony}
\label{sec:ceremony}

A certificate tree proves a completed computational object. It does not, by
itself, show when officials committed to the rule or whether intermediate cuts
were disclosed before later choices could be substituted. We therefore define
a ceremony around the same canonical tree.

Before geographic cuts are released, a genesis transcript commits the rule-
profile hash, root-instance hash, district count, floor/ceiling tree, minimum
review interval, witness threshold, and challenge policy. For a plan with $k$
districts, the completed binary tree contains exactly $k-1$ non-leaf cut
decisions. Round zero releases the root. Round $d$ releases every non-leaf cut
at depth $d$, ordered by binary path, only after the frozen interval has elapsed
since the prior publication.

Every announcement embeds its complete certified split instance, certificate,
and proof. The round also binds its timestamp, external witness receipts, and
the prior round hash. An observer can consequently verify the public prefix,
including exact parent-derived child universes, without receiving unreleased
descendants. The final round changes the transcript to \texttt{completed} and
binds the complete certified-tree identifier. A challenged process may instead
become permanently \texttt{halted}; its accepted prefix and nonempty halt reason
remain verifiable, and resumption requires a new ceremony identifier.

Depth batching is deliberate. Requiring a separate statutory delay for all
$k-1$ cuts would make a large delegation needlessly serial, while disclosing
all cuts at once would eliminate intermediate review. Same-depth cuts depend on
already accepted parents and are mutually independent, so they may be reviewed
together. Different States may likewise proceed in parallel.

The reference implementation exposes immutable \texttt{start},
\texttt{announce}, \texttt{verify}, and \texttt{halt} transitions. It checks
minimum elapsed time and distinct receipt bindings, but opaque receipt strings
are not proof of real-world publication. An official deployment must validate
the signatures and clocks of its named transparency logs or timestamp
authorities.

This ceremony sharpens the Huntington--Hill analogy. The political institution
selects and publishes the governing rule; each subsequent geographic decision
is a locally checkable consequence of that rule. It does not make the rule
value-free, prove a partisan-symmetric outcome, or displace legal and public
review.

For prospective bakeoffs, ceremony evidence forms a separate procedural
scoreboard: proof coverage, prefix verification, release completion, transcript
size, verifier time, halt behavior, and hostile-test resistance. These fields
must not be combined with map-quality metrics into a post hoc scalar winner.
The preexisting bakeoff results predate this protocol and are not retroactively
classified as ceremony runs.
